\section{Lys-task, adaptiv skanning og direkte Pi-test}
\label{sec:lysskann}
Denne oppfølgingen samme dag bygger videre på avsnitt~\ref{sec:piunntak}.
Bruker bekreftet uendret sensor-/servooppsett, frakoblet motorstrøm, at
sensorforsyningen også var av, og at fire LED-er med fire 1\,k$\Omega$-
motstander er tilgjengelige, men ikke montert. Demonstrasjon er oppgitt å
være om omtrent to--tre uker og eksamen om omtrent én måned; eksakte datoer
og gjeldende vurderingskriterier er ikke bekreftet.

\subsection{Egen task og beskyttet lysstatus}
P0 er satt av til to røde baklys, P12 til venstre og P16 til høyre blinklys.
Hver LED får egen seriemotstand; de to røde deler bare signal og jord.
Ingen sensorpinne gjenbrukes. Med antatt $V_{DD}=3{,}3$\,V og
$V_f=1{,}8$\,V blir strømmen per rød LED
\[
 I=\frac{V_{DD}-V_f}{R}=\frac{3{,}3-1{,}8}{1000}=1{,}5\,\mathrm{mA},
\]
altså omtrent 3\,mA samlet fra P0. Beregningen forutsetter de oppgitte
spenningene. Utgangene er konfigurert med nRF52 high-drive H0H1; faktisk
strøm og synlig lys må kontrolleres mot monteringen~\cite{pins}.

\kode{Adaptive_Lights} har en task som alene eier disse tre GPIO-utgangene.
Motor og supervisor publiserer en enumerert tilstand gjennom
\kode{Adaptive_Lamp_Control.State}, et protected object med korte
operasjoner uten I2C, logging eller venting. Dette er en faktisk bruk av
beskyttet delt tilstand~\cite{adaprotected2026}. En feilstatus kan ikke
overskrives av en etterfølgende kjørestatus. En ubehandlet feil ved
taskgrensen oversettes til stoppforespørsel og \kode{Lights_Task_Error}
hos supervisoren; feilen forventes ikke å propagere direkte inn i hovedtasken.

Lys-tasken bruker absolutt \kode{delay until}, ønsket periode 50\,ms og
laveste ordinære taskprioritet. Tapte utgivelser hoppes over. Blinkperioden
er ett sekund med 50\,\% påtid; fasen starter på ved skifte til sving.
Bremselysene viser stoppintensjon, ikke målt retardasjon. Ved latchet feil
lyser begge røde fast mens blinklysene blinker samtidig.
50\,ms er et mykt designmål: sensorpolling og bussarbeid på høyere prioritet
kan forsinke tasken. Dette er ingen ferdig schedulability-analyse av bilen.

\subsection{Begrenset søk fremfor faste skannerunder}
Den aktive \kode{Scan_Path} bruker nå \kode{Adaptive_Quick_Scan}, en generic
med innsprøytet sensorlesing, klokke, stoppflagg og servokontroll. Det gjør
den faktiske søkealgoritmen testbar uten maskinvare. Algoritmen stopper
motorene, kontrollerer sentrerte sensorer og avslutter uten sidepunkter hvis
veien allerede er klar eller svingområdet er utrygt. Ellers prøver den siden
med størst målt sideklaring først. Ved likhet beholdes forrige svingretning.
Kandidatparene er høyre $72^\circ/60^\circ$ og venstre $96^\circ/108^\circ$.

En ugyldig eller blokkert stråle forkaster paret straks. Et bredt nok par
må også bestå en ny sentrert sensorkontroll. Eldste avlesning må være høyst
2000\,ms gammel. Søk avsluttes ved første akseptable kandidat eller etter
høyst fire ekstra sideposisjoner. Tre råavlesninger per posisjon og
sentrerte kontrollmålinger kommer i tillegg. De tidligere aktive faste
5-/7-punktsrundene er dermed erstattet; historiske policytester beholdes.
Servoens tidligere 0,5\,s innstillingstid og 65\,ms globale triggeravstand
er beholdt. Sistnevnte er i samsvar med databladets anbefaling om mer enn
60\,ms målesyklus~\cite{hcsr04}.

Korridormodellen viderefører
$w=2d\sin(\Delta\theta/2)$. For $\Delta\theta=12^\circ$ og krav om
$16+2\cdot5=26$\,cm bredde gir dette
\[
 d\geq\frac{26}{2\sin(6^\circ)}\approx124{,}4\,\mathrm{cm}.
\]
Sparsomme stråler er ikke bevis for at hele mellomrommet er fritt.
2000\,ms er et eksplisitt ferskhetsbudsjett for en stasjonær kandidat,
ikke en garanti for bevegelige hindringer.

Søkekriteriet er kort tid til en tilstrekkelig klarert retning, ikke globalt
maksimal korridorbredde. Hvis to kandidater har kostnad $c_i$ og
suksessannsynlighet $p_i$, blir forventet kostnad
\[
 E_{R\rightarrow L}=c_R+(1-p_R)c_L,\qquad
 E_{L\rightarrow R}=c_L+(1-p_L)c_R.
\]
Høyre først er da gunstig når $p_R/c_R\geq p_L/c_L$. Ingen slike
sannsynligheter er målt; sideklaring brukes foreløpig som en deklarert
heuristikk. En påstand om global matematisk optimalitet ville derfor være
udokumentert.

\subsection{Grunnlag for jevnere kjøring og verifikasjon}
Kontinuerlig kjøring rett frem stopper fortsatt ikke for hver gyldig
front--venstre--front--høyre-kontroll. Den nye ekstra skanningen er
\emph{stasjonær}. Før skanning i bevegelse trengs målt fart, servo- og
sensorlatens og bremsing. Med kalibrert retardasjon $a>0$, reaksjonstid $t$
og fri avstand $d>m$ kan stoppmodellen løses til
\[
 v_{\max}=\sqrt{(at)^2+2a(d-m)}-at.
\]
Dette er et videre designgrunnlag, ikke en aktiv hastighetsregulator.
PWM er ingen fysisk fartsmåling.

Totalt består 2000 native Ada-kontroller og 40 Python-tester. Den nye
lampetesten omfatter modusskifter, blinkgrenser, to samtidige publiserende
tasks og sticky feilstatus. Søkets tester omfatter tidlig avslutning,
ukjente stråler, endret klarering, foreldede/fremtidige tidsstempler,
stopp under søk og propagation av maskinvarefeil. Fire firmwareprosjekter
bygger: hovedprogram, sensordiagnostikk V2, GPIO V3 og lysdiagnostikk.

Via Pi-en ble hele flashen sikkerhetskopiert og GPIO V3s lastesegmenter
kontrollert byte for byte før reset. En full fysisk opptaksserie ga
320/320 \kode{NO_RISE}: front 80, venstre 160 og høyre 80.
Bruker bekreftet at sensorstrømmen var av. Resultatet er dermed en
kommunikasjons-/timeoutkontroll, ikke en sensorfeil, kalibrering eller
bekreftelse på fungerende avstandsmåling. Ingen referanseavstand, MAE eller
RMSE er konstruert. Rådata og kontrollsummer er bevart i repoets evidensmappe.

Separat lysdiagnostikk ble deretter lastet på micro:biten med motor- og
sensorforsyning av. Første flash-forsøk feilet med erase-timeout og er
bevart som mislykket forsøk. Etter eksplisitt reset-and-halt før
programmering lyktes overføringen; full tilbakeavlesning av begge
lastesegmenter samsvarte med ELF-bildet. Den eksakte årsaken til første
feil er ikke endelig fastslått. Ingen automatisk debug-opplåsing eller
helchip-sletting ble brukt.

I 20 sekunder ble 362 GPIO-utgangsregisterprøver samlet sammen med
seriell moduslogg. Alle seks mønstre bestod kontrollen: parkert, venstre,
høyre, kjøring, brems og demonstrert feilblink. LED-ene var ikke montert;
resultatet gjelder utgangslatcher og programkjøring, ikke lysstyrke,
motstandsverdier, elektrisk signalnivå eller fysisk LED-funksjon.
Lysdiagnostikken har ingen motor-/servo-I2C-kommandoer, og står igjen på
kortet ved avslutningen av denne testen. Full bilstyring er ikke flashet.

Selve forelesnings-PDF-en og kilde-ZIP-en er fremdeles ikke hentet:
ChatGPT-nettleseren er innlogget, men prosjektlisten viser lastefeil.
Fagkoblingen må sluttkontrolleres når filene er tilgjengelige lokalt.
